\section{Introduction}
\label{sec:intro}

Much of practical machine learning is classification: approve or reject a loan, route a ticket,
flag a transaction. The discipline that grew around it assumes a model that outputs a probability
for each class. Calibration diagnostics, cost-sensitive decision rules, label-efficient
evaluation and model-agnostic explanations all consume that distribution.

Large language models are now routinely used as zero-shot classifiers, including on tabular data
serialized to text~\citep{hegselmann2023tabllm}. But a language model does not output a class
distribution. It defines a measure over token sequences, and a label is recovered only after
parsing its output. Mass leaks to strings that name no class, and the class distribution can at
best be estimated by sampling or by asking the model to verbalize numbers, which are coarse and
often poorly calibrated~\citep{jiang2021know,tian2023justask,xiong2024uncertainty}; its
self-reported explanations need not reflect what drove the
prediction~\citep{turpin2023unfaithful}.

Some models take a different route: given an input and a predefined set of options, they return a
probability for \emph{every} option~\citep{typesafe_sdk}. Our thesis is that this is not a
cosmetic difference. Moving the output from a measure on strings to a measure on the finite class
space $(C, 2^C)$ is what the classical toolbox requires, including the part of it whose job is to
say when a model should not be trusted.

We contribute (i) a short formal account of the measure a generative model induces on classes and
of the invalid mass it leaks (Section~\ref{sec:background}); (ii) an evaluation of one such
model, Jev, zero-shot on four public datasets, next to an open-weight language model scored so
that it also returns a class distribution and to two supervised models, on identical examples and
with every label budget stated (Sections~\ref{sec:setup}--\ref{sec:metrics}); (iii) a comparison
of macro-F1, cost and latency against the number of training labels, and an ensemble that
escalates low-confidence predictions (Sections~\ref{sec:comparison}--\ref{sec:escalation}); and
(iv) threshold selection, importance-sampled evaluation and perturbation-based explanations
driven only by the class probabilities (Sections~\ref{sec:thresholding}--\ref{sec:explain}).

The tools do not flatter the model. They show that its probabilities are accurate and close to
calibrated on text and biased on the two tabular tasks, in one case in a way a threshold repairs
and in the other not; that a textbook importance-sampling proposal does not pay off on these
probabilities, and why; that a self-hosted open model matches it on text at many times the
price per prediction; and that escalating its uncertain predictions helps exactly where its
confidence is informative. The paper is about which tooling a class distribution makes available
and what that tooling reveals. It is not a leaderboard: two models are zero-shot and two are
supervised.
